\documentclass[11pt,a4paper]{article}

%% ── Encoding & Fonts ──
\usepackage[utf8]{inputenc}
\usepackage[T1]{fontenc}
\usepackage{tgtermes}
\usepackage{newtxmath}
\renewcommand{\ttdefault}{pcr}

%% ── Page Layout ──
\usepackage[margin=1in,headsep=0.4cm]{geometry}
\usepackage{microtype}
\raggedbottom

%% ── Colour ──
\usepackage{xcolor}
\definecolor{navy}{HTML}{1B3A5C}
\definecolor{darknavy}{HTML}{0F2440}
\definecolor{gold}{HTML}{B8860B}
\definecolor{lgray}{HTML}{F5F5F5}
\definecolor{mgray}{HTML}{D0D0D0}
\definecolor{slate}{HTML}{475569}

%% ── Tables ──
\usepackage{booktabs}
\usepackage{tabularx}
\usepackage{caption}
\captionsetup[table]{labelfont={color=navy,bf},textfont=normalfont,skip=4pt}

%% ── Lists ──
\usepackage{enumitem}

%% ── Code Listings ──
\usepackage{listings}

\lstdefinelanguage{json}{
    sensitive=true,
    keywords={},
    morecomment=[l]{//},
    morestring=[b]",
}

\lstset{
    basicstyle=\ttfamily\footnotesize,
    backgroundcolor=\color{lgray},
    frame=single,
    rulecolor=\color{mgray},
    breaklines=true,
    numbers=left,
    numberstyle=\tiny\color{slate},
    numbersep=8pt,
    tabsize=4,
    captionpos=b,
    aboveskip=1.2em,
    belowskip=1.2em,
    xleftmargin=3pt,
    framexleftmargin=3pt,
}

%% ── Section Formatting ──
\usepackage{titlesec}
\titleformat{\section}{\Large\bfseries\color{navy}}{\thesection}{1em}{}[\vspace{-0.15em}\rule{\textwidth}{0.4pt}]
\titleformat{\subsection}{\large\bfseries\color{navy!85}}{\thesubsection}{1em}{}
\titleformat{\subsubsection}{\normalsize\bfseries\color{navy!70}}{\thesubsubsection}{1em}{}

%% ── Headers & Footers ──
\usepackage{fancyhdr}
\setlength{\headheight}{22pt}
\pagestyle{fancy}
\fancyhf{}
\fancyhead[L]{\small\color{slate}\nouppercase{\leftmark}}
\fancyhead[R]{\small\color{slate}\thepage}
\renewcommand{\headrulewidth}{0.4pt}
\renewcommand{\headrule}{\hbox to\headwidth{\color{navy!35}\leaders\hrule height 0.4pt\hfill}}
\fancypagestyle{plain}{%
    \fancyhf{}
    \fancyhead[L]{\small\color{slate}\nouppercase{\leftmark}}
    \fancyhead[R]{\small\color{slate}\thepage}
    \renewcommand{\headrulewidth}{0.4pt}
    \renewcommand{\headrule}{\hbox to\headwidth{\color{navy!35}\leaders\hrule height 0.4pt\hfill}}
}

%% ── Hyperlinks ──
\usepackage{hyperref}
\hypersetup{
    colorlinks=true,
    linkcolor=navy,
    urlcolor=navy,
    citecolor=navy,
    pdftitle={Depression Risk Screening Using Ensemble Machine Learning},
    pdfauthor={Muhammad Ghyas Ali},
    pdfsubject={Machine Learning Project Report -- Mental Health Risk Assessment}
}

%% ── Misc ──
\usepackage{parskip}

%% ── Counters ──
\setcounter{tocdepth}{2}
\setcounter{secnumdepth}{3}

%% ===================================================================
\begin{document}

%% ==========================  COVER PAGE  ============================
\thispagestyle{empty}
\newgeometry{margin=1in}
\begin{center}
\vspace*{3.5cm}

{\color{navy}\rule{0.7\textwidth}{0.8pt}}\\[0.25cm]
{\color{gold}\rule{0.7\textwidth}{1.2pt}}\\[2cm]

{\fontsize{30}{36}\selectfont\bfseries\color{darknavy}
Depression Risk Screening\\[0.2em]
Using Ensemble Machine Learning}

{\color{gold}\rule{0.45\textwidth}{0.5pt}}\\[1.2em]

{\Large\color{darknavy!70} A Production-Grade Web Application\\[0.15em]
for Mental Health Risk Assessment}\\[2.2cm]

{\color{navy}\rule{0.7\textwidth}{1.2pt}}\\[0.25cm]
{\color{gold}\rule{0.7\textwidth}{0.8pt}}

\vfill

\begin{minipage}{0.6\textwidth}
\centering
\renewcommand{\arraystretch}{1.35}
\begin{tabular}{rl}
\textbf{Author}          & Muhammad Ghyas Ali \\[0.3em]
\textbf{Version}         & 3.0 (Optimized Ensemble) \\[0.3em]
\textbf{Date}            & \today \\[0.3em]
\textbf{Source Code}     & \href{https://github.com/MuhammadGhyasAli/ML-Model}{\texttt{github.com/MuhammadGhyasAli/ML-Model}} \\[0.3em]
\textbf{Live Demo}       & \href{https://studentdepressionmodel.vercel.app}{\texttt{studentdepressionmodel.vercel.app}} \\
\end{tabular}
\end{minipage}

\vspace{2.5cm}
\end{center}
\restoregeometry

%% ==========================  ABSTRACT  ==============================
\newpage
\section*{Abstract}
\addcontentsline{toc}{section}{Abstract}
\markboth{ABSTRACT}{ABSTRACT}

\vspace{0.3em}
{\centering\color{gold}\rule{0.35\textwidth}{0.5pt}\par}
\vspace{1.2em}

\noindent
This report presents the design, implementation, and deployment of a machine learning--powered web application for depression risk screening. The system uses a soft-voting ensemble of four classifiers---Logistic Regression, Random Forest, XGBoost, and Gradient Boosting---trained on approximately 27,900 clinical records. A Flask web interface accepts 16 demographic, academic, and lifestyle features, applies 16 engineered feature transformations, and returns a probabilistic risk assessment through an interactive dashboard. The application is deployed on Vercel's serverless infrastructure as \href{https://studentdepressionmodel.vercel.app}{\textbf{ML Model}}, achieving an AUC-ROC of 99.2\%.

\vspace{1.2em}
\noindent
{\small\textbf{Keywords:} Depression screening, ensemble learning, soft voting, XGBoost, Flask, serverless deployment, mental health, feature engineering.}

\vspace*{\fill}

%% ===================  TABLE OF CONTENTS  ============================
\newpage
\tableofcontents
\newpage

%% ====================  LIST OF TABLES  ==============================
\listoftables
\newpage

%% ==========================  INTRODUCTION  ==========================
\section{Introduction}

Mental health disorders, particularly depression, represent a growing global health challenge. Early detection and intervention significantly improve treatment outcomes, yet traditional screening methods remain inaccessible to many due to cost, stigma, and limited clinical resources. This project addresses this gap by developing an accessible, web-based screening tool that leverages ensemble machine learning to provide instant depression risk assessments.

\subsection{Objectives}

\begin{itemize}[leftmargin=*,nosep]
    \item Develop a high-accuracy ensemble model for depression risk classification
    \item Engineer clinically meaningful features from raw demographic and lifestyle data
    \item Build a responsive, production-grade web interface for data collection and result presentation
    \item Deploy the complete system on serverless infrastructure for global accessibility
    \item Provide both a web form and REST API for integration with other applications
\end{itemize}

\subsection{Technology Stack}

\begin{table}[htbp]
\centering
\begin{tabular}{lll}
\toprule
\textbf{Component} & \textbf{Technology} & \textbf{Version} \\
\midrule
Backend Framework & Flask & 3.x \\
Machine Learning & scikit-learn & 1.6.x \\
Gradient Boosting & XGBoost (CPU) & 3.2.x \\
Data Processing & pandas \& numpy & 2.x / 1.24+ \\
Model Serialization & joblib & 1.3+ \\
Frontend & HTML5 / CSS3 / JavaScript & --- \\
Deployment & Vercel (Serverless) & --- \\
Version Control & Git / GitHub & --- \\
\bottomrule
\end{tabular}
\caption{Technology stack used in the project.}
\end{table}

%% ===========================  DATASET  ==============================
\section{Dataset}
\label{sec:dataset}

\subsection{Source and Composition}

The model is trained on the \textbf{Student Depression Dataset}~\cite{depression-dataset} from Kaggle, comprising approximately 27,900 clinical records. Each record contains 16 features spanning three categories:

\begin{itemize}[leftmargin=*,nosep]
    \item \textbf{Demographic:} Gender, Age, City, Profession, Degree
    \item \textbf{Academic and Occupational:} Academic Pressure, Work Pressure, CGPA, Study Satisfaction, Job Satisfaction, Work/Study Hours
    \item \textbf{Lifestyle and Health:} Sleep Duration, Dietary Habits, Financial Stress, Suicidal Ideation, Family History of Mental Illness
\end{itemize}

\subsection{Target Variable}

The target is a binary classification: \texttt{0} (no depression) or \texttt{1} (depression detected). The dataset is balanced through stratified sampling techniques during training.

%% ===================  ML PIPELINE  ==================================
\section{Machine Learning Pipeline}
\label{sec:pipeline}

\subsection{Architecture Overview}

The prediction pipeline is implemented as an \texttt{sklearn.pipeline.Pipeline} consisting of a preprocessing stage followed by a voting ensemble classifier.

\begin{lstlisting}[caption={High-level pipeline structure},language=Python]
Pipeline([
    ('preprocessor', ColumnTransformer([
        ('num', Pipeline([
            SimpleImputer(strategy='median'),
            StandardScaler()
        ]), numeric_features),
        ('ord', Pipeline([
            SimpleImputer(strategy='most_frequent'),
            OrdinalEncoder(handle_unknown='use_encoded_value')
        ]), ordinal_features),
        ('nom', Pipeline([
            SimpleImputer(strategy='most_frequent'),
            OneHotEncoder(handle_unknown='ignore')
        ]), nominal_features)
    ])),
    ('classifier', VotingClassifier(
        estimators=[
            ('lr',  LogisticRegression(C=0.1)),
            ('rf',  RandomForestClassifier(
                n_estimators=400, max_depth=12)),
            ('xgb', XGBClassifier(
                n_estimators=200, max_depth=4,
                learning_rate=0.05)),
            ('gb',  GradientBoostingClassifier(
                n_estimators=300, max_depth=5,
                learning_rate=0.05))
        ],
        voting='soft',
        weights=[1.0, 2.0, 2.0, 1.5]
    ))
])
\end{lstlisting}

\subsection{Preprocessing}

The \texttt{ColumnTransformer} applies distinct preprocessing strategies to three feature groups:

\begin{itemize}[leftmargin=*,nosep]
    \item \textbf{Numeric features} (11 raw + 16 engineered = 27 columns): Missing values are imputed using the median, and features are standardized to zero mean and unit variance.
    \item \textbf{Ordinal feature} (Sleep Duration): Imputed with the most frequent value and encoded ordinally.
    \item \textbf{Nominal features} (City, Profession, Dietary Habits, Degree): Imputed with the most frequent value and one-hot encoded with \texttt{handle\_unknown='ignore'} for robustness against unseen categories.
\end{itemize}

\nopagebreak
\subsection{Ensemble Classifiers}

\subsubsection{Logistic Regression}
A linear model with L2 regularization ($C=0.1$) serving as a baseline. It provides interpretable probability estimates and captures linear relationships between features and depression risk.

\subsubsection{Random Forest}
An ensemble of 400 decision trees with a maximum depth of 12. This model captures non-linear feature interactions and provides feature importance rankings. It receives a weight of 2.0 in the voting ensemble.

\subsubsection{XGBoost}
A gradient-boosted tree model with 200 trees, maximum depth of 4, and a learning rate of 0.05. XGBoost is particularly effective for structured tabular data due to its built-in regularization and optimized gradient boosting algorithm. Its weight in the ensemble is 2.0.

\subsubsection{Gradient Boosting}
A sequential ensemble of 300 trees with a maximum depth of 5 and learning rate of 0.05. This model corrects errors made by previous trees, providing strong predictive performance. Its weight is 1.5.

\subsection{Voting Strategy}

The ensemble uses soft voting, where each classifier outputs class probabilities, and the final prediction is a weighted average:

\begin{equation*}
P(y = 1) = \frac{\displaystyle\sum_{i=1}^{4} w_i \; p_i(y = 1)}{\displaystyle\sum_{i=1}^{4} w_i}
\end{equation*}

where $w_i$ are the classifier weights and $p_i$ are the predicted probabilities. This approach produces calibrated probability estimates.

%% ======================  FEATURE ENGINEERING  ======================
\section{Feature Engineering}
\label{sec:features}

The application includes a dedicated \texttt{engineer\_features()} function that computes 16 derived features from the 16 raw inputs. These engineered features capture clinically meaningful interactions and risk patterns.

\begin{table}[htbp]
\centering
\small
\begin{tabularx}{\textwidth}{lX}
\toprule
\textbf{Feature} & \textbf{Formula / Purpose} \\
\midrule
Total Pressure & Academic Pressure + Work Pressure \\
Pressure--Satisfaction Ratio & Total Pressure / (Study Satisfaction + Job Satisfaction + 1) \\
Sleep Deprived & 1 if Sleep Duration is ``Less than 5 hours'' \\
Sleep Quality & Ordinal encoding of sleep duration (0--3) \\
High Risk & Suicidal thoughts AND Financial Stress $\ge$ 4 \\
Stress Multiplier & Academic Pressure $\times$ Work Pressure \\
Overall Satisfaction & Study Satisfaction + Job Satisfaction \\
Satisfaction--Pressure Ratio & (Satisfaction + 1) / (Total Pressure + 1) \\
Work--Life Balance & Overall Satisfaction / (Work/Study Hours + 1) \\
Financial--Academic Interaction & Financial Stress $\times$ Academic Pressure \\
Suicide Risk Severity & Suicidal thoughts $\times$ (Financial + Total Pressure) / 2 \\
Mental Health Vulnerability & Weighted composite of 5 risk factors \\
CGPA Performance Gap & (4.0 -- CGPA) $\times$ Academic Pressure \\
Perfect Health & All favorable conditions met \\
Extreme Stress & 3+ of 5 extreme conditions present \\
CGPA Achievement & CGPA / (Academic Pressure + 1) \\
\bottomrule
\end{tabularx}
\caption{The 16 engineered features and their computation.}
\end{table}

%% ====================  WEB APP ARCHITECTURE  =======================
\section{Web Application Architecture}
\label{sec:webapp}

\subsection{Framework and Structure}

The web application is built with Flask and follows a standard MVC-like pattern. The entry point is \texttt{app.py}, which defines four page routes and two API endpoints.

\begin{table}[htbp]
\centering
\begin{tabular}{llll}
\toprule
\textbf{Route} & \textbf{Method} & \textbf{Description} & \textbf{Response} \\
\midrule
\texttt{/} & GET & Assessment form & HTML \\
\texttt{/predict} & POST & Form submission, prediction & HTML result \\
\texttt{/about} & GET & Model architecture & HTML \\
\texttt{/documentation} & GET & API reference & HTML \\
\texttt{/api/predict} & POST & JSON prediction API & JSON \\
\texttt{/health} & GET & Health check & JSON \\
\bottomrule
\end{tabular}
\caption{Application routes and their purposes.}
\end{table}

\nopagebreak
\subsection{REST API}

The \texttt{/api/predict} endpoint accepts POST requests with a JSON body containing all 16 input features:

\nopagebreak
\begin{lstlisting}[caption={Example API request and response},language=JSON,numbers=none]
POST /api/predict
{
    "Gender": 1,
    "Age": 22,
    "Academic Pressure": 3,
    "Work Pressure": 2,
    "CGPA": 3.2,
    "Study Satisfaction": 4,
    "Job Satisfaction": 3,
    "Sleep Duration": "7-8 hours",
    "Dietary Habits": "Moderate",
    "Have you ever had suicidal thoughts ?": 0,
    "Work/Study Hours": 8,
    "Financial Stress": 2,
    "Family History of Mental Illness": 0
}

Response 200:
{
    "prediction": 0,
    "depression_probability": 0.1308,
    "no_depression_probability": 0.8692
}
\end{lstlisting}

\subsection{User Interface}

The frontend features a polished enterprise dashboard design with:
\begin{itemize}[leftmargin=*,nosep]
    \item Three-section assessment form with a progress indicator
    \item Real-time client-side field validation
    \item Animated probability bar on the results page
    \item Color-coded risk badges (green, orange, red)
    \item Responsive layout with mobile sidebar navigation
    \item Contextual clinical recommendations based on prediction
    \item Print-friendly result formatting
\end{itemize}

\subsection{Validation Pipeline}

All 16 input fields undergo rigorous validation:
\begin{itemize}[leftmargin=*,nosep]
    \item \textbf{Presence check:} All required fields must be present
    \item \textbf{Type check:} Numeric fields must contain valid numbers
    \item \textbf{Range check:} Values must fall within predefined ranges (e.g., CGPA 0.0--4.0, Age 10--50, Pressure 1--5)
\end{itemize}

%% =========================  DEPLOYMENT  ============================
\section{Deployment}
\label{sec:deployment}

\subsection{Serverless Infrastructure}

The application is deployed on \textbf{Vercel}, a serverless platform, using the Python runtime. The deployment presented several challenges that were addressed through optimization:

\nopagebreak
\subsubsection{Bundle Size Optimization}

The primary deployment challenge was Vercel's 500~MB Lambda ephemeral storage limit. The full \texttt{xgboost} package with GPU support weighs 254~MB (compressed wheel), making the total bundle exceed 990~MB. This was resolved by:
\begin{itemize}[leftmargin=*,nosep]
    \item Using \texttt{xgboost-cpu} (5.6~MB) instead of the full \texttt{xgboost} package
    \item Pinning scikit-learn to version $<$1.7 to maintain pickle compatibility
    \item Resulting function size: \textbf{244~MB} (well under the 500~MB limit)
\end{itemize}

\subsubsection{Version Compatibility}

The model pickle file was trained with scikit-learn 1.6.x, which references the private class \texttt{\_Remainder\-Cols\-List} removed in version 1.7. The scikit-learn version was pinned to $\ge$1.3, $<$1.7 to ensure seamless model loading.

\subsection{Access Links}

\begin{itemize}[leftmargin=*,nosep]
    \item \textbf{Live Application:} \href{https://studentdepressionmodel.vercel.app}{\texttt{studentdepressionmodel.vercel.app}}
    \item \textbf{Source Code:} \href{https://github.com/MuhammadGhyasAli/ML-Model}{\texttt{github.com/MuhammadGhyasAli/ML-Model}}
\end{itemize}

%% ======================  RESULTS & PERFORMANCE  ===================
\section{Results and Performance}
\label{sec:results}

\subsection{Model Metrics}

The ensemble model achieves the following performance metrics:

\begin{table}[htbp]
\centering
\begin{tabular}{lc}
\toprule
\textbf{Metric} & \textbf{Value} \\
\midrule
AUC-ROC & 99.2\% \\
Training Samples & 27,900 \\
Features (Raw + Engineered) & 32 \\
Ensemble Size & 4 classifiers \\
Voting Strategy & Soft (probability-weighted) \\
Validation Method & Stratified K-Fold \\
Model Version & 3.0 (Optimized) \\
\bottomrule
\end{tabular}
\caption{Model performance metrics.}
\end{table}

\subsection{Risk Classification Tiers}

The model's probability output is mapped to three risk levels:

\begin{table}[htbp]
\centering
\begin{tabular}{lc}
\toprule
\textbf{Depression Probability} & \textbf{Risk Level} \\
\midrule
$<$ 50\% & LOW RISK \\
50\%--74\% & MODERATE RISK \\
$\ge$ 75\% & HIGH RISK \\
\bottomrule
\end{tabular}
\caption{Risk classification thresholds.}
\end{table}

%% =======================  PROJECT STRUCTURE  =======================
\section{Project Structure}
\label{sec:structure}

The complete project is organized as follows:

\nopagebreak
\begin{lstlisting}[numbers=none,basicstyle=\ttfamily\footnotesize]
Student Depression/
  app.py                                   # Flask application (287 lines)
  student_depression_ensemble_v3_optimized.pkl  # Trained ML pipeline
  requirements.txt                         # Python dependencies
  vercel.json                              # Vercel deployment config
  wsgi.py                                  # WSGI entry point
  templates/
    index.html                             # Assessment form (3-section)
    result.html                            # Prediction result page
    about.html                             # Model documentation
    documentation.html                     # API reference
\end{lstlisting}

%% ==========================  CONCLUSION  ===========================
\section{Conclusion}
\label{sec:conclusion}

This project successfully demonstrates the development and deployment of a production-grade depression risk screening tool. Key achievements include:

\begin{itemize}[leftmargin=*,nosep]
    \item An ensemble machine learning model achieving 99.2\% AUC-ROC through the combination of four diverse classifiers
    \item A comprehensive feature engineering pipeline that derives 16 clinically meaningful features from raw inputs
    \item A responsive, professionally designed web interface with real-time validation, risk visualization, and contextual recommendations
    \item A REST API for programmatic access, enabling integration with other healthcare platforms
    \item Successful deployment on Vercel's serverless infrastructure, overcoming bundle size constraints through package optimization
\end{itemize}

The tool provides an accessible, evidence-based screening mechanism that can serve as a preliminary step in mental health assessment. It is important to emphasize that this is a screening aid, not a clinical diagnostic instrument. All predictions should be interpreted by qualified mental health professionals.

\nopagebreak
\subsection{Future Work}

\begin{itemize}[leftmargin=*,nosep]
    \item Integration of explainability methods (SHAP/LIME) to provide feature-level insights for each prediction
    \item Expansion of the training dataset to improve generalizability across diverse populations
    \item Addition of longitudinal tracking to monitor changes in risk scores over time
    \item Implementation of multi-language support for broader accessibility
    \item Continuous model retraining pipeline for ongoing performance improvement
\end{itemize}

\vspace{1.5em}
{\centering
\rule{0.35\textwidth}{0.4pt}\\[0.8em]
{\footnotesize\color{slate}
\textit{This project is for screening purposes only and is not a substitute for professional medical advice, diagnosis, or treatment.}}\par}

%% =========================  REFERENCES  ============================
\clearpage
\begin{thebibliography}{99}
\addcontentsline{toc}{section}{References}

\bibitem{depression-dataset}
Hopesb. \textit{Student Depression Dataset}. Kaggle, 2024. \\
\href{https://www.kaggle.com/datasets/hopesb/student-depression-dataset}{\texttt{kaggle.com/datasets/hopesb/student-depression-dataset}}.

\bibitem{source-code}
Muhammad Ghyas Ali. \textit{Depression Risk Screening --- ML Model}. GitHub, 2025. \\
\href{https://github.com/MuhammadGhyasAli/ML-Model}{\texttt{github.com/MuhammadGhyasAli/ML-Model}}.

\bibitem{live-demo}
Muhammad Ghyas Ali. \textit{Depression Risk Screening --- Live Application}. Vercel, 2025. \\
\href{https://studentdepressionmodel.vercel.app}{\texttt{studentdepressionmodel.vercel.app}}.

\end{thebibliography}

\end{document}
